\subsubsection{A scheme with two large extra dimensions confronted with neutrino physics --- arXiv:hep-ph/0208211}

\textbf{Authors:} J. Maalampi, V. Sipil\"ainen, I. Vilja

\paragraph{主な主張}
2次元トーラス上の4枚のブレーンをbulkニュートリノで結ぶ模型では、ニュートリノ混合がブレーン配置や余剰次元の大きさに依存し、当時の実験・宇宙論的制限の下で大きな余剰次元を保つと質量二乗差が非常に小さく制限されることを示した\cite{Maalampi:2002Torus}。

\paragraph{新規性}
5次元の二ブレーン模型をvortexとantivortexが各2個存在する6次元の四ブレーン模型に拡張し、非零KKモードを積分消去してブレーン間混合と2次元KK和の対数的cutoff依存性を含む有効理論を導出した。

\paragraph{位置付け}
2次元bulkニュートリノによるブレーン間混合の初期研究であり、矩形$T^2$のKK和とUV依存性を調べる参考文献となる一方、一般の複素構造や高次KK共鳴の直接観測は扱っていない。
